\documentclass[11pt,a4paper]{article}

% ---- 한글 + 수식 ----
\usepackage{kotex}
\usepackage{amsmath,amssymb,amsfonts,bm}
\usepackage{graphicx}
\usepackage{geometry}
\geometry{margin=25mm}
\usepackage[hidelinks]{hyperref}
\usepackage{enumitem}
\usepackage{tcolorbox}
\usepackage{xcolor}
\usepackage{booktabs}

\renewcommand{\thefootnote}{\arabic{footnote}}

% "한 줄 요약" 박스 (파랑)
\newtcolorbox{readbox}[1][]{
  colback=blue!4, colframe=blue!40!black, boxrule=0.4pt,
  left=6pt,right=6pt,top=4pt,bottom=4pt, #1}
% "주의/반전" 박스 (주황)
\newtcolorbox{notebox}[1][]{
  colback=orange!6, colframe=orange!55!black, boxrule=0.4pt,
  left=6pt,right=6pt,top=4pt,bottom=4pt, #1}
% "충격 발견" 박스 (빨강)
\newtcolorbox{warnbox}[1][]{
  colback=red!5, colframe=red!55!black, boxrule=0.6pt,
  left=6pt,right=6pt,top=4pt,bottom=4pt, #1}

\title{\textbf{SFTF 논문의 ``정직한 한계'' 이야기}\\[4pt]
\large 버그 하나가 어떻게 실험 결과를 바꾸었나\\[2pt]
\normalsize ── 고등학생을 위한, 오늘 발견한 내용 쉽게 읽기 ──}
\author{}
\date{2026-07-04}

\begin{document}
\maketitle

\begin{readbox}
\textbf{이 글의 목표.} 짝꿍 글 ``\emph{지지흐름 텐서장(SFTF) 쉽게 읽기}''에서 우리는
3D 프린팅에서 서포트가 적게 나오는 출력 방향을 \emph{빠르게 추천}하는 방법(SFTF)을
공부했습니다. 그런데 그 논문을 학술지에 낸 뒤, \textbf{계산기 프로그램 안에 숨어 있던
작은 버그} 하나를 발견했습니다. 그 버그를 고치고 실험을 전부 다시 해 보니, 논문의
자랑거리 중 일부가 \emph{사실은 그 버그 때문에 부풀려져 있었다}는 것이 드러났습니다.
이 글은 그 이야기를 \textbf{고등학교 1$\sim$2학년 수학}만으로 읽을 수 있게 풀어 씁니다.
과학에서 ``틀린 것을 정직하게 고치는 일''이 왜 중요한지도 함께 느껴 보면 좋겠습니다.
\end{readbox}

\tableofcontents
\bigskip

%======================================================================
\section{먼저: 무엇이 문제였나 (아주 작은 그릇 이야기)}
%======================================================================

\subsection{컴퓨터가 숫자를 담는 ``그릇''에는 크기가 있다}

컴퓨터는 숫자를 \textbf{정해진 크기의 그릇}에 담습니다. 우리가 쓴 정밀 계산기
프로그램(논문에서 \textbf{TOMO}\footnote{TOMO: 우리가 ``정답''으로 삼은 정밀 서포트
부피 계산기. 모든 방향을 촘촘히(예: 1$^\circ$ 간격) 다 계산해서 진짜 서포트 양을
알려 주지만, 그만큼 느립니다. SFTF는 이 느린 계산을 \emph{어디에 집중할지} 골라 주는
역할입니다.}라고 부릅니다)는, 서포트 부피를 셀 때 중간 계산값을
\textbf{int16}\footnote{int16(16비트 정수): 정수를 담는 작은 그릇의 한 종류.
$-32{,}767$ 부터 $+32{,}767$ 까지만 담을 수 있습니다. 16비트($=2$바이트)로 표현할 수
있는 정수의 범위입니다.}이라는 작은 그릇에 담고 있었습니다.

\begin{notebox}
\textbf{int16 그릇의 한계.} 이 그릇에는 약 $\pm 3$만 까지만 들어갑니다. 그런데
서포트 부피를 셀 때, 물체를 세로로 자른 \emph{기둥(컬럼)} 하나하나에서 높이 값들을
\textbf{계속 더해} 나갑니다. 면(삼각형)이 아주 많은 큰 모델에서는 이 합이 $3$만을
훌쩍 넘겨 버립니다.
\end{notebox}

\subsection{그릇이 넘치면 ``음수''가 튀어나온다 (오버플로)}

시계를 생각해 봅시다. 11시에서 2시간을 더하면 13시가 아니라 \textbf{1시}로 돌아갑니다.
그릇이 넘치면 이렇게 \emph{빙 돌아} 엉뚱한 값이 됩니다. int16도 똑같아서, $32{,}767$ 에서
$1$ 을 더하면 갑자기 $-32{,}768$ 이 됩니다. 이것을 \textbf{오버플로}\footnote{오버플로
(overflow, 넘침): 계산 결과가 그릇(자료형)의 최대치를 넘어, 반대쪽 끝(음수)으로
``휘감겨'' 완전히 틀린 값이 되는 현상. 프로그래밍에서 아주 흔하면서도 찾기 어려운
버그입니다.}라고 합니다.

\begin{warnbox}
\textbf{결정적 증상.} 서포트 부피는 \emph{부피}이므로 절대 음수가 될 수 없습니다.
그런데 면이 100만 개쯤 되는 큰 모델에서 이 계산기가 \textbf{음수 수백만}이라는
말도 안 되는 값을 내놓고 있었습니다. 예를 들어 어떤 큰 부품은 서포트 부피가
$-3{,}105{,}304$ 로 나왔습니다. 부피가 마이너스 삼백만이라니요.
\end{warnbox}

%======================================================================
\section{이 버그가 왜 ``결과 전체''를 흔들었나}
%======================================================================

\subsection{망가진 값들을 자동으로 버리고 있었다}

우리 실험 절차에는 ``\emph{최적값이 $0$ 이하($=$비양수)인 데이터는 이상하니까
빼고 분석한다}''는 안전장치가 있었습니다. 좋은 규칙처럼 보이지만, 문제는
\textbf{그 규칙이 걸러 버린 것들의 정체}였습니다.

\begin{readbox}
\textbf{한 줄 요약.} ``이상해서 버린'' 데이터가 사실은 \emph{버그 때문에 음수가 된
큰 모델들}이었습니다. 즉 우리는 \textbf{큰 기계 부품을 거의 다 빼놓고}, \emph{작은
유기형 모델}(토끼, 용, 인체 스캔 같은 매끈한 곡면들)만으로 채점하고 있었던 셈입니다.
\end{readbox}

작은 유기형 모델들에서는 SFTF가 꽤 잘 작동했습니다. 그래서 ``SFTF가 좋다''는 결론이
나왔지요. 하지만 그건 \textbf{SFTF가 잘하는 문제만 골라서 시험을 본} 것과 비슷했습니다.
버그를 고치자 숨어 있던 큰 부품들이 시험장에 다시 들어왔고, 점수가 확 달라졌습니다.

\subsection{버그를 고치니 일어난 일}

계산기의 그릇을 \textbf{int32}\footnote{int32(32비트 정수): 약 $\pm 21$억까지 담는
훨씬 큰 그릇. 그릇만 키워도 이 문제는 사라집니다.}라는 큰 그릇으로 바꾸었습니다.
그러자:
\begin{itemize}[leftmargin=1.4em]
  \item 큰 부품의 음수들이 \textbf{전부 사라지고} 정상적인 양수 부피가 되었습니다.
  \item 작은 모델들도 값이 \textbf{15$\sim$20\% 정도 달라졌습니다} — 즉 \emph{모든}
        저장된 계산 결과가 조금씩 오염돼 있었던 것입니다.
  \item 이제 큰 부품들이 정식 채점 대상에 들어오면서, 아래처럼 결과가 뒤집힙니다.
\end{itemize}

%======================================================================
\section{뒤집힌 결론 3가지 (제일 중요)}
%======================================================================

\subsection{반전 1 — ``단순한 방법''이 SFTF를 이겼다}

SFTF의 핵심 자랑은 ``\emph{좋은 방향을 똑똑하게 골라 준다}''는 것이었습니다. 이걸
공정하게 시험하려고, \textbf{같은 개수}의 방향만 검사하되 방향을 고르는 방법을 바꿔
비교했습니다.
\begin{itemize}[leftmargin=1.4em]
  \item \textbf{SFTF}: 점수 순위로 상위 20개 방향을 고른다.
  \item \textbf{균등(uniform)}\footnote{균등 방식: 공 위에 점을 되도록 \emph{고르게}
        흩뿌리듯 방향을 잡는 것. 특별한 계산 없이 그냥 ``골고루''입니다.}: 구면에
        방향을 골고루 펼쳐서 고른다.
  \item \textbf{무작위(random)}: 그냥 아무렇게나 고른다.
\end{itemize}

\begin{warnbox}
\textbf{충격적인 결과.} 버그를 고친 뒤 큰 부품까지 포함해 다시 재보니, \textbf{균등
방식이 SFTF를 완전히 눌렀습니다.} (숫자가 작을수록 좋음, 21개 모델 평균)
\begin{center}
\begin{tabular}{lrr}
\toprule
방법 & 평균 비율 & 최악 비율\\
\midrule
SFTF 상위 20 & 3.34 & 34.0\\
\textbf{균등} & \textbf{1.20} & \textbf{2.10}\\
무작위 & 1.32 & 2.56\\
\bottomrule
\end{tabular}
\end{center}
즉 ``\emph{똑똑하게 고른}'' SFTF보다 ``\emph{그냥 골고루}'' 고른 쪽이 더 좋은
방향을 찾았습니다. 큰 기계 부품에서 특히 그랬습니다.
\end{warnbox}

왜일까요? 뒤에서 설명하지만, 핵심은 이렇습니다: SFTF의 점수가 큰 부품에서는
\emph{쓸모가 없어서}, 상위 20개가 한 군데에 \textbf{옹기종기 뭉쳐} 좁은 영역만
들여다봅니다. 반면 균등 방식은 넓게 흩어져 있어서, 정답이 어디 있든 \emph{그물처럼}
걸러 냅니다.

\subsection{반전 2 — 큰 부품에서는 ``특징''이 정답을 못 맞힌다}

SFTF는 각 방향마다 몇 가지 \textbf{특징 점수}(오버행이 얼마나 많은지, 바닥이 얼마나
받쳐야 하는지 등)를 계산해서 순위를 매깁니다. 이 특징이 ``진짜 서포트 양''과 얼마나
잘 맞는지(스피어만 상관\footnote{스피어만 상관: 두 값이 ``같이 커지고 같이
작아지는지''를 등수로 재는 값. $+1$ 이면 완벽히 같은 방향, $0$ 이면 무관, $-1$ 이면
정반대. 즉 $0$ 에 가까우면 ``이 특징으로는 정답을 못 맞힌다''는 뜻입니다.})을
그룹별로 재 보았습니다.

\begin{center}
\small
\begin{tabular}{lcc}
\toprule
특징 & 유기형(토끼) & 큰 기계부품(평균)\\
\midrule
바닥판 항 $B$ (논문의 ``최고 예측자'') & $0.42$ (잘 맞음) & $\approx 0$, 어떤 건 $-0.40$\\
전체 점수 $J$ & $0.62$ (잘 맞음) & $\approx 0$ (못 맞음)\\
\bottomrule
\end{tabular}
\end{center}

\begin{notebox}
\textbf{무슨 뜻인가.} 논문이 ``가장 믿을 만한 예측자''라고 내세운 바닥판 항 $B$ 조차,
큰 기계 부품에서는 정답과 \textbf{거의 관계가 없었습니다}(상관 $\approx 0$). 한 부품
(E10)에서는 심지어 \emph{반대로}($-0.40$) 움직였습니다. 특징이 원래 매끈한
유기형에서만 통했고, \textbf{각진 기계 부품으로는 그 실력이 옮겨 가지 않은} 것입니다.
점수가 엉터리니 순위도 엉터리 → 반전 1의 ``뭉침''이 여기서 나옵니다.
\end{notebox}

\subsection{반전 3 — 제일 어렵다던 모델이 사실은 안 어려웠다}

투고본에서 ``가장 어려운 사례''로 크게 다룬 모델(Happy, 인체 스캔)이 있었습니다.
그 모델의 정답 서포트 부피가 $2{,}112$ 라고 적혀 있었는데, 이 값도 \textbf{음수 오염의
잔재}였습니다. 제대로 고치니 정답은 $8{,}513$ 이었고, SFTF가 처음 고른 방향의 비율은
$8.52$ 배가 아니라 겨우 $1.84$ 배였습니다. 즉 \textbf{원래 그렇게 어려운 문제가
아니었습니다.} 논문의 그림·표에 있던 ``$17{,}991$ 대 $2{,}112$, 비율 $8.52$'' 같은
숫자는 전부 버그가 만든 허상이었습니다.

%======================================================================
\section{그래도 건진 것 (그리고 뜻밖의 발견)}
%======================================================================

\subsection{매끈한 유기형에서는 여전히 쓸모 있다}

다행히 \emph{유기형 모델}(토끼\,·\,용\,·\,인체 같은 매끈한 곡면)에서는 SFTF의 주장이
\textbf{여전히 성립}합니다(숫자만 조금 바뀜):
\begin{itemize}[leftmargin=1.4em]
  \item 바닥판 항 $B$ 가 가장 좋은 단일 예측자 (평균 상관 $0.59$).
  \item 통합 항 $\tilde R=R+B$ 가 가장 강함 ($0.61$).
  \item 상위 20개 방향만 정밀검증하면 격자의 $3.4\%$ 만 보고도 평균 $1.14$ 배 수준을 회복.
  \item 이렇게 학습한 순위가 PCA\,·\,대칭 텐서 같은 기존 방법보다 낫다(단 차이는 줄어듦).
\end{itemize}

\subsection{뜻밖의 발견 — ``버리려던 것''이 더 나았다}

논문의 큰 줄기는 ``\emph{옛날식 대칭 텐서}\footnote{대칭 텐서: 면 법선들의 방향 분포를
요약한 행렬 $M=\sum_i A_i\,m_i m_i^{\!\top}$. 위아래를 뒤집어도 값이 같아서(대칭),
``어느 쪽이 처진 면인지''를 구분 못 한다고 논문은 지적했습니다. 그래서 방향을
구분하는 \emph{비대칭} SFTF로 갈아탄 것이 논문의 출발점입니다.}은 부족하니, 우리의
새로운 비대칭 방법으로 갈아타자''는 것이었습니다. 그런데 큰 부품에서 여러 특징을
시험하다 보니, \textbf{바로 그 버리려던 대칭 텐서가 오히려 정답을 더 잘 맞혔습니다}
(전체 평균 상관 $0.51$ 대 비대칭 $\approx 0$).

\begin{notebox}
\textbf{과학의 반전.} 우리가 ``낡았다''고 대체한 방법이, 정작 어려운 문제에서는 더
쓸모 있었던 것입니다. 물론 대칭 텐서도 완벽하진 않아서(일부 부품에서 실패), 결국
\emph{균등 방식}은 못 이겼습니다. 하지만 ``비대칭이 대칭보다 항상 낫다''던 전제가
\textbf{적어도 기계 부품에서는 틀렸다}는 정직한 발견입니다.
\end{notebox}

%======================================================================
\section{그래서 논문을 어떻게 고쳐야 하나}
%======================================================================

\begin{readbox}
\textbf{한 줄 요약.} SFTF를 ``\emph{모든 부품에서 최고}''가 아니라, \textbf{``매끈한
유기형 부품을 위한, 빠르고 해석 가능한 후보 추천기''}로 \emph{범위를 정직하게 좁히면}
됩니다. 그러면 남은 주장은 모두 사실이고, 방어할 수 있습니다.
\end{readbox}

구체적으로:
\begin{enumerate}[leftmargin=1.5em]
  \item ``SFTF가 단순 방법보다 창을 잘 배치한다''는 \textbf{우위 주장은 철회}한다.
        (큰 부품에서는 균등 방식이 더 낫다는 표를 오히려 정직하게 싣는다.)
  \item 큰 기계 부품에서 특징이 안 통하는 것, 그리고 대칭 텐서가 더 나은 것을
        \textbf{``한계'' 절에 솔직히} 적는다.
  \item 모든 계산 결과를 \textbf{고친 계산기로 다시} 낸 숫자로 교체한다. (완료)
  \item ``어려운 경우 자동 확대(AVE)'' 규칙도 고친 데이터에 맞게 다시 조정한다. (완료)
\end{enumerate}

\begin{warnbox}
\textbf{가장 중요한 교훈.} 결과가 나쁘게 바뀌었다고 숨기지 않는 것. 오히려
``\emph{버그가 우리 방법을 실제보다 좋아 보이게 했었다. 고치고 나니 이런 한계가
드러났다}''라고 밝히는 편이, 훨씬 \textbf{믿을 수 있는 과학}입니다. 심사위원이
언젠가 같은 큰 부품으로 시험했다면 어차피 드러났을 일이고요.
\end{warnbox}

\bigskip
\begin{center}
\emph{---- 끝. ``틀린 것을 정직하게 고치는 일''이 과학의 진짜 실력입니다. ----}
\end{center}

\end{document}
